\chapter*{\centering Project Overview}
\addcontentsline{toc}{chapter}{Project Overview}

Security teams in modern organisations have to make sense of very large volumes of network traffic, alerts and suspicious files. Much of this work is still manual and spread across separate tools, so a single incident can require an analyst to check network monitors, malware scanners and log viewers one after another before the picture becomes clear. When attacks are frequent, or when the activity does not match any known signature, this way of working is slow and error-prone.

The \textit{AI-Based Autonomous Cyber Defense System} proposes a modular platform in which machine learning components analyse different kinds of security evidence and report their findings in a common, structured form. At the mid-term stage, three analysis modules have been implemented and evaluated as working prototypes on public benchmark datasets. The \textbf{Network Threat Classification} module uses a two-stage LightGBM cascade on CIC-IDS2017 flow records to decide whether a flow is benign and, if not, which attack family it belongs to. The \textbf{Malware Classification} module uses CatBoost to separate benign and malicious Windows executables represented by BODMAS feature vectors, followed by a Random Forest that assigns the malware to one of eight categories. The \textbf{Anomaly Detection} module learns only from normal traffic, using an ensemble of unsupervised detectors (Isolation Forest, Local Outlier Factor and an autoencoder) so that it can flag behaviour that has never been seen in training. A log analysis component is under development, while the remaining proposed parts (network monitoring, incident correlation, vulnerability and attack-path analysis, attack prediction, risk analysis, an AI investigation agent and a unified dashboard) are planned for the second half of the project.

The scope of the project is a software prototype developed and tested in a controlled environment with public datasets. It is intended to support, not replace, security analysts. It does not scan production networks, execute countermeasures automatically or develop malware. The modules currently run as separate prototypes with their own command-line or Gradio interfaces; connecting them through a shared backend and dashboard is the main integration goal for the next phase.

The expected impact is a measurable reduction in the manual effort needed to triage security events. Results from the three modules are encouraging. The network classifier reached 99.91\% holdout accuracy with a macro F1-score of 0.958, the malware detector 99.51\% accuracy, and the improved anomaly detector raised its F1-score on deduplicated CIC-IDS2017 traffic from 0.07 to 0.83 while detecting 85\% of attack types never seen during training on NSL-KDD. These results come from benchmark data and still need validation on live traffic, which is part of the future work.
